\section{Pregunta N$^{\circ}$7\qquad Khalid Zaid Izquierdo Ayllón}

\begin{frame}
	\begin{enumerate}\setcounter{enumi}{6}
		\item

		      Un fabricante de bombillas gana \$$0.3$ por cada bombilla
		      que sale de la fábrica, pero pierde \$$0.4$ por cada una
			      que sale defectuosa.
			      Un día en el que fabricó $2100$ bombillas obtuvo un
		      beneficio de \$$484.4$.
		      Determine el número de bombillas buenas y defectuosa
		      según el siguiente requerimiento.

		      \begin{enumerate}[a)]
			      \item

			            Modele el problema.

			      \item
			            Determine la norma matricial de $A$.

			      \item

			            Determine el número de condicionamiento de $A$.

			      \item

			            Resolver usando el método de Factorización
			            $LDL^{T}$ y Cholesky.
		      \end{enumerate}

	\end{enumerate}

	\begin{solution}
		  \begin{enumerate}[a)]
			      \item
			            Modele el problema.

                        $x_{1}$:bombillas buenas \qquad $x_{2}$:bombillas defectuosas
                        
                        \begin{align*}      
                        
                            x_{1}+x_{2}=2100 \qquad\qquad\quad \longrightarrow \qquad x_{1}+x_{2}=2100

                        0.3x_{1}-0.4x_{2}=484.4\qquad \longrightarrow \qquad 3x_{1}+4x_{2}=4844
                        \end{align*}
                        \begin{equation*}
                            \underbrace{\begin{pmatrix}
                                1 & 1\\
                                3 & -4
                            \end{pmatrix}}_{A}
                            \underbrace{\begin{pmatrix}
                                x_{1}\\
                                x_{2}
                            \end{pmatrix}}_{x}
                            =
                            \underbrace{\begin{pmatrix}
                                2100\\
                                4844
                            \end{pmatrix}}_{b}
                        \end{equation*}

                    \item
			            Determine la norma matricial de $A$.

                        $\|A\|_{\infty}=max\{2;7\}=7$
            \end{enumerate}
	\end{solution}
\end{frame}

\begin{frame}
	\begin{solution}
		  \begin{enumerate}[c)]
		    \item 
                    Determine el número de condicionamiento de $A$.

                    Calculamos la inversa de A:

                    \begin{equation*}
                            A^{-1}=
                            \begin{pmatrix}
                                \frac{4}{7} & \frac{1}{7}\\
                                \frac{3}{7} & -\frac{1}{7}
                            \end{pmatrix}
                    \end{equation*}
                    $\|A^{-1}\|_{\infty}=max\{\frac{5}{7};\frac{4}{7}\}=\frac{5}{7}$

                    $Cond(A)=\|A\|\cdot\|A^{-1}\|=7\times \frac{5}{7}=5$               
		  \end{enumerate}
            \begin{enumerate}[d)]
                \item 
                    Resolver usando el método de Factorización
			            $LDL^{T}$ y Cholesky.

                    Dado que la matriz A no es simétrica, realizaremos una operación para que lo sea.

                    \begin{equation*}
                            \underbrace{\begin{pmatrix}
                                1 & 3\\
                                1 & -4
                            \end{pmatrix}}_{A^t}
                            \cdot
                            \underbrace{\begin{pmatrix}
                                1 & 1\\
                                3 & -4
                            \end{pmatrix}}_{A}
                            \underbrace{\begin{pmatrix}
                                x_{1}\\
                                x_{2}
                            \end{pmatrix}}_{x}
                            =\underbrace{\begin{pmatrix}
                                1 & 3\\
                                1 & -4
                            \end{pmatrix}}_{A^t}
                            \cdot
                            \underbrace{\begin{pmatrix}
                                2100\\
                                4844
                            \end{pmatrix}}_{b}
                        \end{equation*}
                        
                        \begin{equation*}
                            \underbrace{\begin{pmatrix}
                                10 & -11\\
                                -11 & 17
                            \end{pmatrix}}_{A'}
                            \underbrace{\begin{pmatrix}
                                x_{1}\\
                                x_{2}
                            \end{pmatrix}}_{x}
                            =
                            \underbrace{\begin{pmatrix}
                                16632\\
                                -17276
                            \end{pmatrix}}_{b'}
                        \end{equation*}
                    
            \end{enumerate}
	\end{solution}
\end{frame}
\begin{frame}
    \begin{solution}
        \begin{enumerate}[d)]
            \item Realizaremos la factorización mediante el uso de un algoritmo aplicado en python, el cu
        \end{enumerate}
    \end{solution}    
\end{frame}